\section{Introduction}

Density functional theory (DFT) and its linear-response
time-dependent extension (LR-TDDFT) provide efficient routes to
ground-state and excitation energies in molecular electronic
structure.\cite{HohenbergKohn1964,KohnSham1965,RungeGross1984,Casida1995}
Their predictive accuracy is ultimately controlled by the
exchange-correlation (XC) approximation, which must provide not
only an energy but also the derivatives that enter the
self-consistent potential and the response kernel. In Kohn--Sham
DFT, the first functional derivative of the XC energy determines
the XC contribution to the SCF potential; in adiabatic LR-TDDFT,
the second derivative enters the Casida response matrices that
govern excitation energies and oscillator strengths. The XC
functional therefore determines the ground-state SCF reference that
underlies subsequent excited-state response calculations.

Finding a suitable XC functional is therefore
central to reliable ground-state SCF and excited-state response
calculations. Neural network XC functionals provide a flexible route
to learn such functional forms, including nonlocal XC structure,
from accurate reference data.\cite{GLOBE,Schmidt2019NonlocalXCML}
For such excited-state targets, however, fitting ground-state
energies alone is not enough: the functional must also produce
stable self-consistent densities, useful frontier orbital gaps, and
a response kernel consistent with the learned XC energy. This
requirement naturally leads to differentiable training, in which the
functional is optimized through the electronic-structure calculation
that will ultimately use it.\cite{Li2021KohnShamRegularizer} A recent
overview of differentiable DFT emphasizes that this changes the
role of the KS equations: the functional is not only evaluated inside
a self-consistent DFT calculation, but trained through the
differentiable SCF map, with gradients from observables propagated
back to the model parameters while the governing electronic-structure
equations remain in the loop.\cite{vonStrachwitz2026DifferentiableDFT}
For TDDFT, the corresponding differentiable program must extend
beyond the ground-state SCF fixed point to include response-kernel
construction, TDA/Casida eigenvalue problems, excitation energies,
oscillator strengths, and spectrum-level objectives. Differentiable
training is therefore important because it can place losses on the
quantities actually used to judge excited-state predictions, rather
than only on pre-response ground-state energies.

This training view is supported by a broader set of differentiable
algorithms in molecular simulation, semiempirical electronic
structure, quantum chemistry, and materials calculations.\cite{JAXMD,DXTB,Athavale2026DifferentiableSemiempirical,PySCFAD,JRystal,JRystalDirectOptimization,Schmitz2025ADDFPT}
For the present problem, the most relevant progress lies in
differentiable DFT and XC functional construction. DQC and D4FT expose
differentiable DFT workflows for learning or optimizing density
functionals,\cite{DQC,D4FT} GradDFT provides a ground-state framework
for machine-learning enhanced density functionals,\cite{GradDFT} and
jax-xc makes conventional Libxc energy-density forms available as
differentiable functional channels.\cite{JAXXC} In parallel, neural XC
models have explored more flexible functional forms, including
nonlocal and coefficient-basis descriptions beyond a fixed semilocal
ansatz.\cite{Schmidt2019NonlocalXCML,GLOBE,KovacsDLDH2026} Together
these works show how XC functionals can be represented, evaluated, and
trained inside self-consistent DFT calculations. The missing piece for
excited states is an end-to-end differentiable TDDFT formulation in
which the neural XC energy, SCF potential, response kernel, and
excited-state observables belong to the same differentiable
computational graph.

We present GradTDDFT to address this gap by extending GradDFT from
ground-state neural XC functional modeling to differentiable
linear-response TDDFT training. The neural XC model is written as a
coefficient-basis functional with trainable semilocal, HF, and
optional MP2/PT2 channels. Its semilocal basis can be built from
jax-xc energy-density components, so conventional LDA, GGA, and
meta-GGA forms can serve as differentiable functional channels inside
the neural model. GradTDDFT also implements implicit differentiation
for self-consistent ground-state training and excited-state response
training, allowing training objectives to be propagated back to the
XC parameters through the SCF and TDDFT maps. In this way, the
learned XC energy, SCF potential, response kernel, and excited-state
observables are treated within one differentiable training framework.
